% Source: Weinberg GR, physical PDF pp. 245--248
%         (printed pp. 222--225).
% Coverage: Section 9.3, The Energy-Momentum Tensor; equations
%           (9.3.1)--(9.3.12).
% Figures/tables/footnotes: reference marker 1; no figures or tables.
% Status: source-reviewed and compile-clean.
% Uncertainties: none.

\subsection{The Energy-Momentum Tensor}\label{sec:9.3}

In order to complete the computational program outlined in
Section~\ref{sec:9.1}, we must show how to calculate the
energy-momentum tensor \(T^{\mu\nu}\) that serves as the source of the
gravitational fields. We shall first consider how the conservation laws of
energy and momentum appear in the post-Newtonian approximation. In general,
the conservation laws read \(\nabla_\mu T^{\mu\nu}=0\), or in more detail
\begin{equation}
  \partial_\mu T^{\mu\nu}
  =
  -\Gamma^\nu{}_{\mu\lambda}T^{\mu\lambda}
  -\Gamma^\mu{}_{\mu\lambda}T^{\lambda\nu}.
  \tag{9.3.1}\label{eq:9.3.1}
\end{equation}
The term of order \(\bar M\bar v/\bar r^4\) with \(\nu=0\) gives
\begin{equation}
  \partial_t T^{00(0)}+\partial_iT^{i0(1)}=0,
  \tag{9.3.2}\label{eq:9.3.2}
\end{equation}
since all connection coefficients are at least of order
\(\bar v^2/\bar r\). This may be regarded as the law of conservation of
mass; it should not surprise us to find mass conserved in the
post-Newtonian approximation, for a large rate of conversion of mass into
energy would produce temperatures at which the particles of the system
moved relativistically, in conflict with the assumption
\(\bar v\ll1\). Apart from its intrinsic importance,
Eq.~\eqref{eq:9.3.2} is indispensable for consistency of the harmonic
coordinate conditions. Equations~\eqref{eq:9.1.53} and
\eqref{eq:9.1.55} imply
\[
  0
  =
  \nabla^2\left(
    -2\partial_tg_{00}^{(2)}+\partial_i g_{0\ii}^{(3)}
  \right)
  =
  \nabla^2\left(
    4\partial_t\phi+\boldsymbol\nabla\cdot\boldsymbol\zeta
  \right).
\]
Since \(\phi\) and \(\boldsymbol\zeta\) must vanish at infinity, this
verifies Eq.~\eqref{eq:9.1.66}.

Returning to Eq.~\eqref{eq:9.3.1}, the \(\nu=\ii\) term of order
\(\bar M\bar v^2/\bar r^4\) gives
\[
  \partial_tT^{0\ii(1)}+\partial_jT^{ij(2)}
  =-\Gamma^{i(2)}{}_{00}T^{00(0)},
\]
or, using Eq.~\eqref{eq:9.1.67},
\begin{equation}
  \partial_tT^{0\ii(1)}+\partial_jT^{ij(2)}
  =
  -(\partial_i\phi)T^{00(0)}.
  \tag{9.3.3}\label{eq:9.3.3}
\end{equation}
Since \(T^{ij}\) is the flux of momentum, this is the conservation of
momentum; the right-hand side is just the Newtonian gravitational force
density, equal to the mass density \(T^{00(0)}\) times
\(-\boldsymbol\nabla\phi\).

There are no other conservation laws involving only the terms in
\(T^{\mu\nu}\) needed to calculate the fields in the post-Newtonian
approximation, namely \(T^{00(0)},T^{00(2)},T^{i0(1)}\), and
\(T^{ij(2)}\). Furthermore, Eqs.~\eqref{eq:9.3.2} and
\eqref{eq:9.3.3} involve the metric only through \(\phi\), which can be
calculated in the Newtonian approximation. Hence the procedure to be
followed is essentially iterative. We first solve the Newtonian equations of
motion, use the solution (plus the equations of state) to determine these
terms in \(T^{\mu\nu}\), compute the post-Newtonian fields \(\psi\) and
\(\boldsymbol\zeta\), recompute the motions of the particles, and so on. It
can be shown\hyperref[ch9-ref-1]{\textsuperscript{1}} that this procedure can
be kept going: to compute the fields in the \(N\)th approximation, we need
terms in \(T^{\mu\nu}\) satisfying conservation laws that involve the
fields only in the \((N-1)\)th approximation.

For completeness, the conservation laws governing the next-highest terms
are, after use of Eqs.~\eqref{eq:9.1.67}--\eqref{eq:9.1.72},
\begin{equation}
  \partial_tT^{00(2)}+\partial_iT^{i0(3)}
  =
  T^{00(0)}\partial_t\phi,
  \tag{9.3.4}\label{eq:9.3.4}
\end{equation}
and
\begin{equation}
  \begin{split}
  \partial_tT^{0\ii(3)}+\partial_jT^{ij(4)}
  ={}&
  -T^{00(0)}
   \left[\partial_i(2\phi^2+\psi)+\partial_t\zeta_i\right]
  -T^{00(2)}\partial_i\phi\\
  &-T^{0j(1)}
   \left[
     \partial_j\zeta_i-\partial_i\zeta_j
     -4\delta_{ij}\partial_t\phi
   \right]\\
  &-T^{jk(2)}
   \left[
     \delta_{jk}\partial_i\phi
     -4\delta_{ik}\partial_j\phi
   \right].
  \end{split}
  \tag{9.3.5}\label{eq:9.3.5}
\end{equation}
Thus the source terms \(T^{00(2)},T^{i0(3)}\), and \(T^{ij(4)}\), which
would be needed to calculate the post-post-Newtonian fields, obey
conservation laws involving the metric only to post-Newtonian order.

We still need a model with which to calculate the energy-momentum tensor.
The simplest such model is an assembly of freely falling particles that
interact only gravitationally and, perhaps, in localized collisions. From
Eq.~\eqref{eq:5.3.5} we have
\begin{equation}
  T^{\mu\nu}(\mathbf x,t)
  =
  \frac1{\sqrt{-g}}
  \sum_n m_n
  \frac{\dd x_n^\mu}{\dd t}
  \frac{\dd x_n^\nu}{\dd t}
  \left(\frac{\dd\tau_n}{\dd t}\right)^{-1}
  \delta^3\!\bigl(\mathbf x-\mathbf x_n(t)\bigr),
  \tag{9.3.6}\label{eq:9.3.6}
\end{equation}
where \(m_n,x_n^\mu(t)\), and \(\tau_n\) are the mass, spacetime position,
and proper time of the \(n\)th particle. Writing
\(-g=1+\mathfrak g^{(2)}+\mathfrak g^{(4)}+\cdots\), a calculation using
Eq.~\eqref{eq:4.7.5} gives
\begin{equation}
  \mathfrak g^{(2)}
  =
  \eta^{\mu\nu}g_{\mu\nu}^{(2)}
  =
  -g_{00}^{(2)}+g_{ii}^{(2)}
  =
  -4\phi.
  \tag{9.3.7}\label{eq:9.3.7}
\end{equation}
Using Eqs.~\eqref{eq:9.3.7} and \eqref{eq:9.2.3} in
Eq.~\eqref{eq:9.3.6}, we find
\begin{equation}
  T^{00(0)}
  =
  \sum_n m_n\delta^3(\mathbf x-\mathbf x_n),
  \tag{9.3.8}\label{eq:9.3.8}
\end{equation}
\begin{equation}
  T^{00(2)}
  =
  \sum_n m_n
  \left(\phi+\frac12v_n^2\right)
  \delta^3(\mathbf x-\mathbf x_n),
  \tag{9.3.9}\label{eq:9.3.9}
\end{equation}
\begin{equation}
  T^{i0(1)}
  =
  \sum_n m_n v_n^i\delta^3(\mathbf x-\mathbf x_n),
  \tag{9.3.10}\label{eq:9.3.10}
\end{equation}
\begin{equation}
  T^{ij(2)}
  =
  \sum_n m_n v_n^iv_n^j\delta^3(\mathbf x-\mathbf x_n),
  \tag{9.3.11}\label{eq:9.3.11}
\end{equation}
where \(\mathbf v_n\equiv\dd\mathbf x_n/\dd t\).

To impose the conservation laws, recall that
% NOTATION EXCEPTION: The full derivative with respect to x_n^i is retained
% because it must be distinguished here from differentiation with respect to
% the field coordinate x^i.
\[
  \partial_t\delta^3(\mathbf x-\mathbf x_n(t))
  =
  v_n^i\frac{\partial}{\partial x_n^i}
  \delta^3(\mathbf x-\mathbf x_n(t))
  =
  -\mathbf v_n\cdot\boldsymbol\nabla
  \delta^3(\mathbf x-\mathbf x_n(t)).
\]
Consequently,
\[
  \partial_tT^{00(0)}+\partial_iT^{i0(1)}=0,
\]
and
\[
  \partial_tT^{0\ii(1)}+\partial_jT^{ij(2)}
  =
  \sum_n m_n\frac{\dd v_n^i}{\dd t}
  \delta^3(\mathbf x-\mathbf x_n).
\]
Thus the mass-conservation equation \eqref{eq:9.3.2} is automatically
satisfied, while Eq.~\eqref{eq:9.3.3} is satisfied if and only if each
particle obeys the Newtonian equation of motion
\begin{equation}
  \frac{\dd\mathbf v_n}{\dd t}
  =
  -\boldsymbol\nabla\phi(\mathbf x_n).
  \tag{9.3.12}\label{eq:9.3.12}
\end{equation}

The program for calculating the motion of a set of gravitating point
particles is therefore:
\begin{enumerate}
  \item[(A)] Solve the Newtonian problem; that is, solve
  Eqs.~\eqref{eq:9.3.12} and \eqref{eq:9.1.58} for
  \(\phi(\mathbf x)\) and \(\mathbf x_n(t)\). (This is the only step that
  is not always straightforward.)
  \item[(B)] Use the results of (A) and
  Eqs.~\eqref{eq:9.3.8}--\eqref{eq:9.3.11} to compute the required terms
  in the energy-momentum tensor.
  \item[(C)] Use the results of (A) and (B), and
  Eqs.~\eqref{eq:9.1.62} and \eqref{eq:9.1.65}, to compute the
  post-Newtonian fields \(\boldsymbol\zeta\) and \(\psi\).
  \item[(D)] Use the results of (A) and (C), and Eq.~\eqref{eq:9.2.1}, to
  calculate the post-Newtonian corrections to the trajectories
  \(\mathbf x_n(t)\).
  \item[(E)] And so on.
\end{enumerate}
